\documentclass[10pt,a4paper]{amsart}
\usepackage[margin=19mm]{geometry}
\usepackage{amsmath,amssymb,mathtools}
\usepackage[scr=rsfs,cal=euler]{mathalfa}
\usepackage{xfrac}
\usepackage[unicode,hidelinks,bookmarks=false]{hyperref}
\newcommand{\ie}{i.e.\ }
\newcommand{\cf}{cf.\ }
\newcommand{\resp}{respectively\ }
\newcommand{\Subsection}{\subsection}
\providecommand{\nobreakdash}{\nobreak\hbox{-}}
\setlength{\emergencystretch}{2em}
\multlinegap0pt
\usepackage{amssymb}
\newtheorem{lem}{Lemma}
\title{Computable Scott Sentences and the Friedman-Stanley embedding}
\author{David Gonzalez, Julia Knight}
\date{Source: arXiv:2605.04404v1 (CC BY 4.0)}
\begin{document}
\maketitle

\noindent\emph{Source and licence.} Computable Scott Sentences and the Friedman-Stanley embedding. Version arXiv:2605.04404v1 (CC BY 4.0), distributed by arXiv under the Creative Commons Attribution 4.0 International licence, http://creativecommons.org/licenses/by/4.0/, as recorded in the arXiv licence metadata for this exact version and shown on the abstract page https://arxiv.org/abs/2605.04404v1. Full licence notice: \texttt{licenses/arxiv-2605.04404-CC-BY-4.0.txt}.\par\smallskip

\noindent\textit{Editorial scope.} Counted unit: the article's lem lem:property'sAxioms (source0.tex lines 542--546). The article's complete original proof of it is delivered with it (source0.tex lines 548--587, one \texttt{proof} environment of 40 lines). No mathematics is written, repaired or re-derived: the statement and the proof are the article's own lines, in the article's own notation, and no retained line is substituted or re-flowed.  No further source-local context is needed: every cross-reference of every retained line resolves inside the two delivered spans.  Nothing was omitted from any retained span, so the package carries no omission record.  No retained line contains a citation, so the wrapper carries no bibliography and no bibliography entry.  No retained line contains a figure environment, an image-inclusion command or drawing code, so no image is delivered and the package declares no asset dependency.  Editorial formatting only, in addition: an AMS \texttt{amsart} preamble that restates the article's own declarations and macros used by the retained lines, namely the environments \texttt{lem} on separate counters, together with the standard packages those lines need.  The retained environments print in this document as Lemma 1 in order of appearance, whereas the article prints the same items with the numbering of its own full document.  The complete native source of the article as distributed in the arXiv e-print for this exact version is bundled unchanged as \texttt{sources/199859/source0.tex}.
\begin{lem}\label{lem:property'sAxioms}

For a computable ordinal $\alpha\geq 2$, there is a computable $\Pi_\alpha$-sentence stating the properties of Length, Consistency, Replication, plus $\beta$-agreement, and $\beta$-permutation for $\beta < \alpha$.  

\end{lem}

\begin{proof}

\noindent
\textbf{Length}:  Recall that $L$ is a finite relational language.  For each $n$, there are only finitely many possible formulas $D$ assigning an atomic type to the tuple $\bar{x} = (x_1,\ldots,x_n)$.  

\noindent
(1)  First, taking into account all $n$, we say that each $x$ satisfies $U_D$ for exactly one $D$.  This is computable $\Pi_2$.  To say that $x = \lambda$ (the base node), we write $p(x) = x$.  To say that $x$ is at level $n\geq 1$, 
we say $p^{n+1}(x) = p^n(x) \ \&\ p^n(x) \not= p^{n-1}(x)$.  This is finitary quantifier-free.           

\noindent
(2)  We say that for each $n$, for all $x$ at level $n$, $x$ satisfies $U_D$ for one of the finitely many $D$  assigning an atomic type to $(x_1,\ldots,x_n)$.  This is computable~$\Pi_2$.

\noindent
Altogether, it is computable $\Pi_2$ to say that $T$ satisfies (1) and (2).      

\bigskip
\noindent
\textbf{Consistency}:  Recall that $T$ has the Consistency Property if for $\sigma\subseteq\tau$, the atomic type $D'$ assigned by $\tau$ implies the atomic type $D$ assigned by $\sigma$.  For a given $n$, let $P_n$ be the set of pairs $(D,D')$ such that $D$ is an atomic type in the first $n$ variables, $D'$ is an atomic type in the first $n+1$ variables, and $D'$ logically implies $D$.  We say for all $x,y$ if $y$ is at level $n+1$ and $p(y) = x$, then $\bigvee_{(D,D')\in P_n} (U_Dx\ \&\ U_{D'}y)$.  This is computable $\Pi_2$.    

\bigskip
\noindent
\textbf{Replication}:  Recall that $T$ has the Replication Property if for all $\sigma$ if some successor satisfies $U_D$, then infinitely many do.  This is computable $\Pi_2$.

\bigskip
\noindent
\textbf{$\beta$-agreement}:  Recall that $T$ has the $\beta$-agreement Property if for all $n$, for all computable $\Sigma_\beta$ formulas $\psi(\bar{x})$ in variables among the first $n$, and all $\sigma$ at level $n$, if some $\tau\supseteq\sigma$ has $\tau\Vdash\psi(\bar{x})$, then no $\tau'\supseteq\sigma$ forces $neg(\psi(\bar{x}))$. 
It is computable $\Sigma_\beta$ to say of $\tau$ that it extends $\sigma$ and forces $\psi(\bar{x})$.  It is computable $\Pi_{\beta}$ to say of $\tau\supseteq\sigma$ that it forces $neg(\psi(\bar{x}))$.  To say that $\tau$ does not force $neg(\psi(\bar{x}))$ is computable $\Sigma_\beta$.  We have a computable $\Pi_{\beta+1}$ sentence saying for all $n$, for all $\psi(\bar{x})$ with variables among the first $n$ and $\sigma$ at level $n$, if some extension of $\sigma$ forces $\psi(\bar{x})$, then no extension forces $neg(\psi(\bar{x}))$.    

\bigskip
\noindent
 If $\alpha = \beta+1$, then it is enough to say that $T$ has the $\beta$-agreement property, which is $\Pi_{\beta+1}$.  If $\alpha$ is a computable limit ordinal, then we say that $T$ has the $\beta$-agreement property for all $\beta < \alpha$.  This is computable $\Pi_\alpha$.

\bigskip
\noindent
\textbf{$\beta$-permutation}:  Recall that $T$ has the $\beta$-permutation property if for $\psi$ a computable $\Sigma_\beta$ formula, and $\bar{x}$ consisting of the first $n$ variables, if $\sigma$, at level $n$, has an extension forcing $\psi(\bar{x},z)$, where $z$ is a later variable, then some extension of $\sigma$ forces $\psi(\bar{x},x')$, where  $x'$ is the $(n+1)^{st}$ variable.  For fixed $n$, $\psi$ and $\sigma$ at level $n$, it is computable $\Sigma_\beta$ to say that some extension of $\sigma$ forces $\psi(\bar{x},z)$ for some later variable $z$, and it is computable $\Sigma_\beta$ to say that some immediate extension forces $\psi(\bar{x},x')$.  Then it is computable $\Pi_{\beta+1}$ to say that for all $n$, $\psi$, and $\sigma$ at level $n$, if some extension of $\sigma$ forces $\psi(\bar{x},z)$ for some $z$, then some immediate extension of $\sigma$ forces $\psi(\bar{x},x')$.

\bigskip
\noindent
 If $\alpha = \beta+1$, it is enough to say that $T$ has the $\beta$-permutation property, which is computable $\Pi_{\beta+1}$.  If $\alpha$ is a computable limit ordinal, then we say that $T$ has the $\beta$-permutation property for all $\beta < \alpha$.  This is computable $\Pi_\alpha$.          
\end{proof}

\end{document}
